\section{Study Design}\label{sec:design}
This section describes the fixed components, the documentation treatments, the tasks and oracles, and the statistical plan (Fig.~\ref{fig:overview}, left). We ran two studies. \textbf{Study~1} (exploratory; 114 synthetic tasks, 2{,}052 trials) surfaced two confounds in our own treatments and an uptake failure in on-demand documentation. \textbf{Study~2}, which the rest of this paper reports, was pre-registered after Study~1. It fixes those confounds, adds real-world tasks derived from KernelBench, and turns the Study~1 post-hoc conditions into planned ones. The protocol, file hashes and all trial logs are released.

\subsection{Fixed IR, toolchain and agent loop}
All conditions share one toolchain and one agent loop. Programs are StableHLO modules with a single \texttt{@main}. We parse and verify them with the MLIR/StableHLO bindings in jaxlib 0.11.2 (StableHLO 1.20.0) and execute them on the XLA compiler's CPU backend (XLA:CPU) in a sandboxed subprocess with a 60\,s timeout. An operation allow-list admits StableHLO and three element-wise operations from CHLO, the companion ``client'' dialect, and excludes \texttt{custom\_call}. The task fixes the signature exactly. The agent receives a system prompt, the documentation (if any) and the task, and replies with a program. Failures are returned verbatim (the verifier diagnostic, or the first mismatching oracle element), with up to two repairs. Documentation stays in context, and in the interactive conditions the retrieval tool remains available for up to 6~calls. Each trial is a fresh conversation at temperature~0, one per task and condition.

\subsection{Canonical facts and interface treatments}
The treatments differ only in how one fact inventory reaches the agent. The inventory has 96 facts over 23 operation groups (25 syntax, 16 type, 32 shape, 23 semantics). Each fact is anchored to a clause of the pinned specification~\cite{stablehlospec} and authored twice: as a natural-language sentence and as a formal predicate in specification notation (Fig.~\ref{fig:overview}, right). Table~\ref{tab:conds} lists the eight treatments.

Two changes remove the Study~1 confounds. First, every syntax fact now carries the same general schema and the same concrete instance in both renderings (e.g.\ both \texttt{array<i64: ...>} and \texttt{array<i64: 1>}); in Study~1, prose carried only the instance and the formal rendering only the schema. Second, the structured contract is compact JSON, so its mean length is $0.98\times$ that of prose over all 328 tasks (Study~1: $1.31\times$), measured with the o200k tokenizer. An automated audit checks every task: identical fact-ID sequences, each rendering present verbatim, and no IDs leaking into prose. All 29 code examples verify with the pinned toolchain. A task's fact set is the cross-cutting \texttt{module} facts (program structure, types, the integer-list attribute encoding, no implicit broadcasting) plus the facts of the operation groups the task lists as relevant. Every condition receives the identical task text, signature and list of relevant groups.

\begin{table}[t]
\caption{Interface treatments (Study~2). Tok.: mean documentation tokens (o200k tokenizer) over 328 tasks.}
\label{tab:conds}
\centering\footnotesize
\setlength{\tabcolsep}{3pt}
\begin{tabular}{@{}llllr@{}}
\toprule
& Condition & Facts upfront & Rendering & Tok. \\
\midrule
-- & \cond{None} & none & -- & 0 \\
A & \cond{Minimal} & syntax & prose & 951 \\
N & \cond{Syntax-only str.} & syntax & JSON & 850 \\
B & \cond{Prose} & all & prose & 1455 \\
C & \cond{Structured} & all & JSON + IDs & 1469 \\
D & \cond{Interactive} & none (on demand) & JSON, tool & 0 \\
D$^\ast$ & \cond{Int.\ forced} & as D, 1st call forced & JSON, tool & 0 \\
D$^+$ & \cond{Int.\ core} & module; rest on demand & JSON, tool & 362 \\
\bottomrule
\end{tabular}
\end{table}

\subsection{Tasks and independent oracles}\label{sec:tasks}
The evaluation population combines controlled synthetic tasks with real-world tasks lowered from PyTorch.

\textbf{Synthetic (114 tasks).} We wrote 38 parametric templates over broadcasting, reductions, contractions, layout, slicing and padding, masking and indexing, conversion, and composites (e.g.\ softmax). Each template is instantiated three times with distinct dimension sizes, and one additional instance per template forms a development split. Each task is specified mathematically. Its oracle is a numpy reference evaluated on three held-out cases that include semantic traps: all-negative rows for max-reductions, exact zeros for strict comparisons, and out-of-range \texttt{dynamic\_slice} starts.

\textbf{KernelBench-derived (153 tasks).} We converted all 200 level-1 and level-2 tasks of KernelBench~\cite{ouyang2025kernelbench} (MIT license). They cover single operators, normalizations, pooling, convolutions, losses, and fused operator chains. The pipeline (i)~shrinks integer size literals deterministically, preserving divisibility where required; (ii)~exposes all module parameters as explicit \texttt{@main} arguments, keeping the training-mode semantics under which KernelBench evaluates; (iii)~uses PyTorch eager outputs on three cases, each with freshly initialized parameters and inputs, as the oracle ($\mathrm{rtol}{=}10^{-3}$); and (iv)~produces a reference StableHLO program via \texttt{torch.export} and \texttt{torchax}. The reference is never shown to agents; it only demonstrates that the task is expressible and that the oracle and toolchain agree. The agent sees the PyTorch module source, the argument mapping and the signature. We excluded 24 tasks, each with a logged reason: 18 were too large even after shrinking, 2 had shape couplings broken by shrinking, 2 use dropout, 1 takes a non-tensor input, and 1 failed validation because its reference disagreed with PyTorch. Of the remaining 176, a 10\% hash split (23 tasks) forms the development set, leaving 153 evaluation tasks and a total evaluation population of 267 tasks.

\textbf{Validation (gate G1).} All 328 reference programs (152 synthetic, 176 KernelBench) parse, verify, execute and match their oracles. We also measured oracle sensitivity by mutating the reference programs (operator swaps, constant and attribute perturbations). On synthetic tasks the oracle rejected all 132 non-equivalent executing mutants; the 13 survivors are semantically equivalent. On KernelBench tasks it rejected 322 of 432 executing mutants. Of the 107 distinct survivors, 90 matched the reference on five fresh signed inputs (likely equivalent, e.g.\ dead running-statistics updates) and 7 could not be probed. The remaining 10, in 9 tasks, were missed because KernelBench draws inputs from \texttt{torch.rand}, which is non-negative. Re-grading the three agent programs accepted on those tasks with signed inputs flipped one (mini, Minimal condition), a single false acceptance among 5{,}000 analyzed first attempts. We report the pre-registered grading.

\subsection{Models, budget and statistics}
We used three dated snapshots of one model family, \texttt{gpt-4.1-nano/-mini/-2025-04-14}, called \emph{nano}, \emph{mini} and \emph{4.1} below. A local Qwen2.5-7B solved 0 of 12 development trials and was not used. The total API budget was capped at \$35 and enforced by the runner. Trials were ordered task-blocked (tasks shuffled, conditions shuffled within each task), so an outcome-blind budget stop leaves complete, paired tasks. The primary endpoint is first-attempt semantic correctness, and the primary contrast is C$-$B. Eight paired contrasts were pre-specified (Table~\ref{tab:contrasts}). For each model we report the mean paired difference in percentage points (pp) with a 95\% cluster-bootstrap CI (10k resamples; clusters are templates for synthetic tasks and individual tasks for KernelBench), together with exact McNemar tests, Holm-corrected within each model. Per-suite results are secondary and Holm-corrected within each suite. The failure taxonomy (assembly syntax, attribute encoding, literal, shape/type, unknown op, wrong values) was fixed before the runs.
